\section{A subpower bound for the triangle trace}
\label{trace:section}

For $p<\pc$ write
\[
 T=T_p,\qquad m_p=\norm{T_p},\qquad a_j=\tr(T^j).
\]
The quantity $a_3$ is the type average of the triangle diagram.
We will prove that it grows more slowly than any prescribed
positive power of $(\pc-p)^{-1}$.

Triangle diagrams and positive connection matrices are classical components of the
tree-graph method of Aizenman and Newman \cite{AizenmanNewman1984};
quantitative relations between triangle diagrams and susceptibility
are developed in \cite{Hutchcroft2022Triangle}. The closed-contact
estimate below is related to the exploration-score and critical-tail
refinements of the two-ghost method
\cite{Hutchcroft2020Locality,Hutchcroft2021PowerLaw}. We prove the
specific contact and trace bounds used here.

For a nonnegative invariant function $F$ of an ordered pair of
distinct finite clusters, define the positive summation functional
\[
 \mathcal D(F)=\left[\E_p\frac1{s_o}
                         \sum_{B\ne C_o}F(C_o,B)\right].
\]
The sum ranges over the other open clusters in the configuration.
For a nonnegative invariant kernel $J$, put
\[
 k_J(A,B)=\sum_{x\in A}\sum_{y\in B}J(x,y),
 \qquad k_e=k_{\Adj}.
\]
In particular $k_e$ counts graph edges between the two clusters.
Mass transport gives the useful rerooting formula
\begin{equation}
 \mathcal D(k_JF)=
 \left[\E_p\sum_v J(o,v)F(C_o,C_v)
                         \one_{\{C_o\ne C_v\}}\right].
 \label{trace:rerooting}
\end{equation}
Indeed send from $x$ to each $y\in C_x$ the mass
$s_x^{-1}\sum_vJ(y,v)F(C_x,C_v)\one_{\{C_x\ne C_v\}}$.
The outgoing and incoming sides of \eqref{eq:s02} are precisely
the two sides of \eqref{trace:rerooting}.  All terms are
nonnegative, so no integrability assumption is needed for this
identity.  Another application, sending
$F(A,B)/(\abs A\abs B)$ from every vertex of $A$ to every vertex
of $B$, shows that $\mathcal D(F(A,B))=\mathcal D(F(B,A))$.

\begin{lemma}[Contacts of two large clusters]
\label{trace:contact-tail}
Fix $0<\alpha<1/2$ and $\epsilon_1>0$ such that
$r=1/2+\epsilon_1-2\alpha<0$.  Uniformly for
$\pc/2\le p<\pc$ and $L\ge1$,
\begin{equation}
 \mathcal D\left(k_e^2
            \one_{\{\min(\abs A,\abs B)\ge L\}}\right)
       \le C L^{1/2+\epsilon_1-2\alpha}.
 \label{eq:s24}
\end{equation}
\end{lemma}

\begin{proof}
Fix an oriented graph edge $(x,v)$ and $n\ge1$.  Let
\[
 H_n=\{C_x\ne C_v,\ n\le\min(s_x,s_v)<2n\}.
\]
On distinct clusters the two size events have disjoint finite open
witnesses.  The BK inequality, coupling with critical percolation,
and \eqref{eq:s21} give
\begin{equation}
 \Prob_p(H_n)\le C n^{-2\alpha}.
 \label{trace:two-large-clusters}
\end{equation}
We also need to control their number of contacts, and for this we
compare two exploration scores.

First explore $C_v$ in the full graph, querying every edge incident
to each discovered vertex exactly once, including edges whose
endpoints have both already been discovered.  Stop after
$N=\lceil2dn\rceil$ queries if the exploration has not finished.
Let $S_{\mathrm{full}}$ be the sum of $\omega_f-p$ over the tested
edges.  Every active query is a fresh Bernoulli($p$) bit.  Padding
after stopping by zero increments, the elementary centered
Bernoulli exponential bound gives
\[
 \E_p e^{\lambda S_{\mathrm{full}}}\le e^{N\lambda^2/8},
 \qquad
 \Prob_p(\abs{S_{\mathrm{full}}}>t)
                       \le2e^{-2t^2/N}.
\]
This is an unconditional estimate, with no cluster-size event
imposed.

On the same percolation configuration, perform a second exploration
by first exposing $A=C_x$ together with its internal bits and
closed boundary.  If $v\notin A$, explore the cluster of $v$ in
$V\setminus A$, again querying every incident edge of that domain
and using the same deterministic query cap $N$.  Let the resulting
score be $S_{\mathrm{ext}}$, and set it to zero if $v\in A$.
Conditional on the exposure, all the queried exterior bits are
fresh and independent.  The identical tail bound therefore holds
conditionally, uniformly in $A$, and hence also unconditionally
after averaging over $A$.

On $H_n\cap\{s_v\le s_x\}$, both explorations finish the same
cluster $B=C_v$, because $\abs B<2n$ and its number of incident
edges is at most $d\abs B<N$.  The completed queried sets are,
respectively, all full-graph edges incident to $B$ and all such
edges with both ends outside $A$.  Their difference consists
exactly of the $A$--$B$ contacts, which are closed.  Query order is
irrelevant to these completed sums, so
\[
 S_{\mathrm{full}}-S_{\mathrm{ext}}=-p k_e(A,B).
\]
It follows that the event under consideration with
$k_e(A,B)>2t/p$ is contained in the union of the two score-tail
events.  Take $t=n^{1/2+\epsilon_1}$, and use $p\ge\pc/2$.
Interchanging $x$ and $v$ handles the opposite size order.  Thus
\[
 \Prob_p\left(H_n\cap
             \{k_e(C_x,C_v)>C n^{1/2+\epsilon_1}\}\right)
             \le C e^{-c n^{2\epsilon_1}}.
\]
Each score was concentrated in its own fresh law before
intersecting with $H_n$; no independence between the scores is
asserted.  Since $k_e\le d\min(s_x,s_v)<2dn$ on $H_n$,
\eqref{trace:two-large-clusters} now yields
\[
 \E_p[k_e(C_x,C_v)\one_{H_n}]
 \le C n^{1/2+\epsilon_1-2\alpha}
                      +C n e^{-c n^{2\epsilon_1}}
 \le C n^r.
\]
Apply \eqref{trace:rerooting} with one contact factor and partition
$\min(\abs A,\abs B)\ge L$ into the bins
$[2^jL,2^{j+1}L)$, $j\ge0$.  There are at most $d$ neighbors of
each root, and $\sum_{j\ge0}(2^jL)^r\le C_rL^r$ because $r<0$.
This proves \eqref{eq:s24}.
\end{proof}

\begin{proposition}[Subpower triangle growth]
\label{trace:subpower}
For every $\eta>0$ there is $C_\eta<\infty$ such that
\begin{equation}
 a_3(p)\le C_\eta(\pc-p)^{-\eta},
 \qquad \pc/2\le p<\pc.
 \label{eq:s23}
\end{equation}
\end{proposition}

\begin{proof}
All operators below are bounded at each subcritical parameter.
The operator-norm local absolute continuity supplied by
\eqref{eq:s04}, and cyclicity of the finite trace, give
$a_3'=3\tr(T^2T')$ almost everywhere.  Set $J=T^2$.
For fixed terminals $x,y$, the first identity in \eqref{eq:s04}
can also be written
\[
 (1-p)T'_p(x,y)
   =\E_p[k_e(C_x,C_y)\one_{\{C_x\ne C_y\}}].
\]
Each edge between the two terminal clusters has exactly one
orientation from the first cluster to the second.  Applying
\eqref{trace:rerooting} to $J$ therefore gives
\begin{equation}
 (1-p)a_3'=3\mathcal D(k_e k_J)
 \quad\text{for almost every }p<\pc.
 \label{eq:s25}
\end{equation}

We will use the following square bound:
\begin{equation}
 \mathcal D(k_J^2)\le\tr(JTJT)=a_6\le m_p^3a_3.
 \label{trace:kernel-square}
\end{equation}
To check the first inequality, reroot one $J$ factor using
\eqref{trace:rerooting} and expand the other one.  This gives
\[
 \left[\sum_{v,x',y'}J(o,v)J(x',y')
   \Prob_p(o\leftrightarrow x',\ v\leftrightarrow y',\
                                      C_o\ne C_v)\right].
\]
The two connections on this event have disjoint open witnesses,
so BK bounds their probability by $T(o,x')T(v,y')$.
Symmetry of $J$ and $T$ identifies the resulting sum with
$\tr(JTJT)$.  Because $J=T^2$, this is $a_6$.  Finally, positivity
of $T$ gives $T^6\le m_p^3T^3$ in quadratic-form order, which
proves the remaining inequality after taking the trace.

Choose $\alpha$ and $\epsilon_1$ as in
\cref{trace:contact-tail}.  Split \eqref{eq:s25} according to
whether the smaller terminal cluster has size less than $L$.
For the part where $A$ is small, reroot one contact and expose
$A=C_o$ with its internal bits and boundary.  If $v\notin A$,
its exterior cluster is fresh, so
\[
 \E_p[k_J(A,C_v)\mid A\text{ and its exposure}]
 \le\sum_{y\in A}\sum_z J(y,z)T(v,z)
 =\sum_{y\in A}(T^3)(y,v).
\]
The positive operator $T^3$ satisfies
\[
 \abs{(T^3)(y,v)}
 \le\sqrt{(T^3)(y,y)(T^3)(v,v)}\le c_\sigma a_3,
\]
since every diagonal entry is bounded by $a_3/\min_i\sigma_i$.
Consequently
\[
 \mathcal D(k_e k_J\one_{\{\abs A<L\}})
 \le d c_\sigma a_3\,[\E_p(s_o\one_{\{s_o<L\}})]
 \le C L^{1-\alpha}a_3.
\]
Here $s\one_{\{s<L\}}\le L^{1-\alpha}s^\alpha$ and
\eqref{eq:s21} supply the last bound.  Reversing the ordered
cluster pair gives the same estimate when $B$ is small, since
$k_e$ and $k_J$ are symmetric in their two arguments.

For the remaining part, Cauchy--Schwarz for the positive measure
underlying $\mathcal D$, followed by \eqref{eq:s24} and
\eqref{trace:kernel-square}, gives
\[
 \mathcal D(k_e k_J\one_{\{\min(\abs A,\abs B)\ge L\}})
 \le C L^{1/4+\epsilon_1/2-\alpha}m_p^{3/2}\sqrt{a_3}.
\]
The factor $(1-p)^{-1}$ in \eqref{eq:s25} is uniformly bounded
on $[\pc/2,\pc)$.  Also $a_3\ge1$, directly from the nonnegative
kernel sum and $T(o,o)=1$.  Thus $b=\sqrt{a_3}\ge1$ is locally
absolutely continuous and satisfies
\begin{equation}
 b'\le C L^{1-\alpha}b
        +C L^{1/4+\epsilon_1/2-\alpha}m_p^{3/2}
 \quad\text{almost everywhere}.
 \label{eq:s26}
\end{equation}

We make the exponent choices explicit.  Given $\eta>0$, choose
$0<\zeta<\min(1,\eta)$, and then choose
$1/2-\alpha>0$, $\epsilon_1>0$, and $\eta'>0$ sufficiently small
that $1/2+\epsilon_1-2\alpha<0$ and
\begin{align*}
 A_*&=(2-\zeta)(1-\alpha)<1,\\
 B_*&=(2-\zeta)(1/4+\epsilon_1/2-\alpha)
                       +\tfrac32(1+\eta')<1+\eta/2.
\end{align*}
These choices are possible by continuity: at the limiting values
$\alpha=1/2$, $\epsilon_1=\eta'=0$, the two expressions are
$1-\zeta/2$ and $1+\zeta/4$, respectively.  Put
$\delta=\pc-p$ and $L=\delta^{-(2-\zeta)}\ge1$.
The susceptibility bound \eqref{eq:s22} implies
$m_p\le\chi_{\max}(p)\le C\delta^{-1-\eta'}$, so
\eqref{eq:s26} becomes
\[
 b'(p)\le C\delta^{-A_*}b(p)+C\delta^{-B_*}.
\]
Since $A_*<1$, the integral of the first coefficient from
$\pc/2$ to $\pc$ is finite.  Multiplication by its integrating
factor therefore gives
\[
 b(p)\le C\left(b(\pc/2)+
                 \int_{\pc/2}^p(\pc-r)^{-B_*}\,dr\right)
       \le C_\eta\delta^{-\eta/2}.
\]
For the last inequality one may bound the integrand by a constant
times $(\pc-r)^{-1-\eta/2}$, using $B_*<1+\eta/2$.
Squaring proves \eqref{eq:s23}.
\end{proof}
